\title{Simplified model of the NiC2 planes in LaNiC2}
\author{Henry Sheehy}
\date{ \today }
  
\maketitle
  
\tableofcontents
\newpage

\section{Model}
Jorge described a simplified model of the $\text{NiC}_2$ planes in $\text{LaNiC}_2$, which is depicted in Figure \ref{fig:unit_cell}.
In $k$-space the model is described by the following tightbinding Hamiltonian
\begin{equation}
    \mathcal{H}= -\mu\mathbb{1}_2+\mathcal{T},
\end{equation}
where $\mu$ is the chemical potential, $\mathbb{1}_2$ is the $2\times2$ identity matrix and $\mathcal{T}$ is the hopping matrix
\begin{equation}
    \mathcal{T}= \alpha \sigma_+ +\beta \sigma_- +\text{h.c.},
\end{equation}
$\sigma_\pm$ are the Pauli raising and lowering matrices, $\alpha$ and $\beta$ parameterise the hopping
\begin{eqnarray}
    \alpha &=& -v-w e^{-ik_x}-2t_d\cos{(k_y)}e^{-i k_x},\\
    \beta &=& -2t_d\cos{(k_y)},
\end{eqnarray}
where $v$ is the intradimer hopping, $w$ is the interdimer hopping and $t_d$ is a diagonal hopping which stacks the SSH chains.
In particular, note that the SSH model is topological only when $w>v$, however, in the model Jorge describes, $v>w$.

\createfigure{LaNiC2/unit_cell}{unit_cell}

\newpage
\section{Bandstructure}
The eigenvalues of the Hamiltonian are
\begin{equation}
    E_\pm = \pm\sqrt{\alpha^*+\beta} \sqrt{\alpha+\beta^*},
\end{equation}
which are plotted in Figure \ref{fig:topological_bandstructure} for the set of parameters
$(\mu, v, w, t_d) = (0, 0.6, 1.2, 0.9)$,
i.e.
the SSH topological phase.
In the trivial phase, 
$(\mu, v, w, t_d) = (0, 1.2, 0.6, 0.9)$,
the band structure looks as Figure \ref{fig:trivial_bandstructure}.
At $\mu=0$, the bands are seen to cross at zero-energy, possibly  Dirac node.

% \createfigure{LaNiC2/topological_bandstructure}{topological_bandstructure}
% \createfigure{LaNiC2/trivial_bandstructure}{trivial_bandstructure}

\createfigure{LaNiC2/bandstructure_phase_transition}{bandstructure_phase_transition}

\newpage
\section{Finite size results: ubiquitous topological modes}
Performing a finite size calculation with $41\times41$ sites, with periodic boundary conditions perpendicular to the axis of the SSH dimers, and open-boundary conditions along their axis, we expected to observe topological features, namely Majorana and Bogoliubov-Weyl modes,
due to the simple model's simiarities to a recently studied model [CITE: Rosenberg and Manousakis arXiv:2203.12004v1].
However, in addition to similar modes obeserved in both models' topological phases 
(see Figure \ref{fig:k_space_topological}),
in our simple model
we surprisingly observe the same topological features in the parameter regime corresponding to the trival phase of the SSH model,
see Figure \ref{fig:k_space_trivial}.

\createfigure{LaNiC2/k_space_topological}{k_space_topological}

\createfigure{LaNiC2/k_space_trivial}{k_space_trivial}

Crossing the phase transition, we observe no change in the bandstructure Figure \ref{fig:bandstructure_phase_transition}, so the transition appears to be topological.

\createfigure{LaNiC2/phase_diagram_closed}{phase_diagram_closed}

\createfigure{LaNiC2/phase_diagram_open}{phase_diagram_open}
